% !TEX program = xelatex
% Bowen Lv — Curriculum Vitae (English) — compile with xelatex
\documentclass[10pt,a4paper]{article}

%==================== Fonts ====================
\usepackage{fontspec}
\setmainfont{texgyretermes}[
  Extension = .otf, UprightFont = *-regular, BoldFont = *-bold,
  ItalicFont = *-italic, BoldItalicFont = *-bolditalic]
\setsansfont{texgyreheros}[
  Extension = .otf, UprightFont = *-regular, BoldFont = *-bold,
  ItalicFont = *-italic, BoldItalicFont = *-bolditalic]
% Chinese fallback for the name only (macOS system fonts)
\usepackage{xeCJK}
\setCJKmainfont{Songti SC}[BoldFont={Heiti SC}]

%==================== Layout ====================
\usepackage[a4paper,top=0.85cm,bottom=0.8cm,left=1.5cm,right=1.5cm]{geometry}
\usepackage[dvipsnames,table]{xcolor}
\usepackage{titlesec}
\usepackage{enumitem}
\usepackage{graphicx}
\usepackage{tikz}
\usepackage{fontawesome5}
\usepackage[hidelinks]{hyperref}
\usepackage{ragged2e}

\definecolor{accent}{RGB}{31,78,214}
\definecolor{accentd}{RGB}{18,46,120}
\definecolor{inkc}{RGB}{27,36,51}
\definecolor{mutedc}{RGB}{91,107,130}
\definecolor{rulec}{RGB}{205,221,248}

\color{inkc}
\hypersetup{colorlinks=true, urlcolor=accent, linkcolor=accent}

\setlength{\parindent}{0pt}
\linespread{1.02}

%==================== Section headings ====================
\titleformat{\section}
  {\sffamily\large\bfseries\color{accentd}}
  {}{0pt}{}
  [{\vspace{2pt}\color{rulec}\titlerule[1.2pt]}]
\titlespacing*{\section}{0pt}{5pt}{3pt}

%==================== Helpers ====================
\newcommand{\entry}[2]{%
  \par\vspace{1pt}%
  {\sffamily\bfseries\color{inkc}#1}\hfill{\small\itshape\color{mutedc}#2}\par
  \vspace{0.5pt}}
\newcommand{\role}[1]{{\small\color{accentd}\sffamily #1}\par\vspace{1pt}}

\setlist[itemize]{leftmargin=1.4em, topsep=1pt, itemsep=1pt, parsep=0pt, label={\color{accent}\small$\bullet$}}

\newcommand{\circphoto}[4]{%
  \begin{tikzpicture}
    \clip (0,0) circle (#2);
    \node at (0,#4) {\includegraphics[width=#3]{#1}};
    \draw[rulec, line width=1.2pt] (0,0) circle (#2);
  \end{tikzpicture}}

\begin{document}
\pagestyle{empty}

%==================== Header ====================
\begin{minipage}[c]{0.72\textwidth}
  {\sffamily\fontsize{26}{30}\selectfont\bfseries\color{accentd}Bowen Lv}\quad
  {\sffamily\large\color{mutedc}吕博文}\\[4pt]
  {\sffamily\color{inkc}Ph.D. Student, Department of Computer Science, Tsinghua University}\\[2pt]
  {\small\color{mutedc}Research: LLM Post-Training\,·\,Reinforcement Learning\,·\,GUI / Computer-Use Agents}\\[8pt]
  {\small
   \faEnvelope[regular]~\href{mailto:lvbowen1228@gmail.com}{lvbowen1228@gmail.com}\quad
   \faPhone*~(+86) 183-3728-3380\\[3pt]
   \faGlobe~\href{https://extreme1228.github.io/extreme1228/}{extreme1228.github.io/extreme1228}\quad
   \faGoogle~\href{https://scholar.google.com/citations?user=5JW6fqwAAAAJ&hl=zh-CN}{Google Scholar}}
\end{minipage}\hfill
\begin{minipage}[c]{0.24\textwidth}
  \raggedleft\circphoto{photo.jpg}{1.85cm}{3.7cm}{-0.55cm}
\end{minipage}

\vspace{2pt}

%==================== Education ====================
\section{Education}
\entry{Tsinghua University \quad Department of Computer Science and Technology (KEG Lab)}{2025 -- Present}
\role{Ph.D. in Computer Science\;·\;Advisor: Prof. Yuxiao Dong\;·\;Focus: LLM + Reinforcement Learning}

\entry{Tongji University \quad School of Electronic and Information Engineering}{2021 -- 2025}
\role{B.Eng. in Computer Science (Elite CS Class)\;·\;Rank 1 / 20 in major}

%==================== Experience ====================
\section{Experience}
\entry{Zhipu AI (智谱) \quad Research Intern --- LLM Post-Training \& GUI Agent RL}{Feb 2025 -- Present}
\role{One year (and ongoing) on large-model post-training and RL for GUI / computer-use agents}
\begin{itemize}
  \item Built core agentic-RL training infrastructure: \textbf{AgentRL}, a from-scratch, highly asynchronous, train--inference-decoupled online RL framework with Dockerized multi-environment deployment and high-concurrency rollout for multi-task, multi-turn training.
  \item Led verifiable task synthesis and online RL for computer-use agents, reaching open-source SOTA on OSWorld / ScienceBoard (\textbf{ScaleCUA}, first author; details below).
  \item Leading \textbf{Auto-Post-Train (APT)} --- an \textbf{AI-for-AI, self-improving} multi-agent post-training system where a Conductor orchestrates env / task / rollout / train workers to fully automate the entire loop --- \textbf{data collection, environment adaptation, task construction, rollout, quality auditing, and model training} --- so models keep self-improving with minimal human intervention.
\end{itemize}

%==================== Publications ====================
\section{Selected Publications}
% —— ScaleCUA highlighted ——
\entry{ScaleCUA: Scaling Computer Use Agents with Verifiable Task Synthesis and Efficient Online RL}{arXiv, 2026}
\role{\textbf{Bowen Lv} (first author), Xiao Liu, Yanyu Ren, Hanyu Lai, Bohao Jing, Hanchen Zhang, et al.\quad\faExternalLink*~\href{https://arxiv.org/abs/2607.11185}{arXiv:2607.11185}}
\begin{itemize}
  \item \textbf{ScaleCUA} pairs verifiable task synthesis with efficient online RL to break the data-scarcity and RL-inefficiency bottlenecks of computer-use agents; its \textbf{VeriGen} pipeline scales to 100+ concurrent workers via a shared Docker interaction probe, yielding 24K+ verifiable tasks and nearly 3K high-quality RL tasks.
  \item Training-side \textbf{Frontier Sampling} (rollouts to the model's learning frontier) and \textbf{Visual Context Segmentation} (sliding-window visual context, 2.83$\times$ speedup); with Qwen3.5-9B, reaches \textbf{68.7\% on OSWorld} and \textbf{54.0\% on ScienceBoard} --- new open-source SOTA, surpassing larger models such as EvoCUA-32B and Kimi K2.5.
\end{itemize}

\entry{AgentRL: Scaling Agentic Reinforcement Learning with a Multi-Turn, Multi-Task Framework}{Preprint, 2025}
\role{Hanchen Zhang, Xiao Liu, \textbf{Bowen Lv}, et al.\quad\faExternalLink*~\href{https://arxiv.org/abs/2510.04206}{arXiv:2510.04206}}
\begin{itemize}
  \item A multi-task, multi-turn agentic-RL framework; after RL, average performance on AgentBench surpasses leading closed-source APIs (Claude-4, GPT-5).
\end{itemize}

\entry{KARL: Reinforcement Learning for LLM Agents on Multi-Turn Knowledge-Intensive Agentic Tasks}{ACL 2026 Long Papers}
\role{Xueqiao Sun, Xiao Liu, \textbf{Bowen Lv}, et al.\quad\faExternalLink*~\href{https://aclanthology.org/2026.acl-long.2196/}{ACL Anthology}}
\begin{itemize}
  \item An RL method for multi-turn, knowledge-intensive agentic tasks; published in ACL 2026 Long Papers.
\end{itemize}

\entry{Focus on What Matters: Separated Models for Visual-Based RL Generalization}{NeurIPS 2024}
\role{Di Zhang, \textbf{Bowen Lv}, et al.\quad\faExternalLink*~\href{https://arxiv.org/abs/2410.10834}{arXiv:2410.10834}}
\begin{itemize}
  \item Separated models that improve generalization of visual RL in complex environments; published at NeurIPS 2024 (\textbf{second author}).
\end{itemize}

\entry{GLM-5V-Turbo: Toward a Native Foundation Model for Multimodal Agents}{Preprint, 2026}
\role{Wenyi Hong, Xiaotao Gu, \ldots, \textbf{Bowen Lv}, \ldots, GLM-V Team\quad\faExternalLink*~\href{https://arxiv.org/abs/2604.26752}{arXiv:2604.26752}}
\begin{itemize}
  \item GLM-V team technical report on a native multimodal-agent foundation model; contributed to GUI-agent post-training.
\end{itemize}

%==================== Awards ====================
\section{Awards \& Competitions}
\entry{Programming Contests}{2022 -- 2024}
\begin{itemize}
  \item CCPC National Invitational (Guangdong) \textbf{Gold}; ICPC Nanjing \textbf{Silver} and Hong Kong \textbf{Silver}; CCCC individual \textbf{National Second Prize} (all 2024).
\end{itemize}
\entry{Honors}{2021 -- 2025}
\begin{itemize}
  \item Shanghai Outstanding Graduate (2025); National Scholarship (2022); Tongji SEIE Qidi Scholarship (2023).
\end{itemize}

\end{document}
